% Status: history (G1b). Superseded by notes/dsl_updates.md + castlegen/channels. Kept for the record.
\documentclass[11pt]{article}
\usepackage[margin=1in]{geometry}
\usepackage{amsmath,amssymb}
\usepackage{booktabs}
\usepackage{dsfont}
\usepackage{algorithm}
\usepackage{algpseudocode}

\newcommand{\Z}{\mathbb{Z}}
\newcommand{\R}{\mathbb{R}}
\newcommand{\1}{\mathds{1}}

\title{Coarse-to-fine sampling: target, generator, sampler}
\author{Hierwave notes}
\date{}

\begin{document}
\maketitle

\noindent Three objects, kept separate throughout:
\begin{enumerate}
  \item the \textbf{target}, an undirected energy model on tiles that says what we want to sample;
  \item the \textbf{generator}, a chain graph over coarse-to-fine latent variables that approximates it;
  \item the \textbf{sampler}, the prolong--perturb--smooth pass that approximately samples the generator.
\end{enumerate}
Code: \texttt{castlegen/legacy/hier.py}, \texttt{texsyn.py} (\texttt{CoordVar}), \texttt{promise.py},
\texttt{quantities/}. Companion to \texttt{texture\_sampling.md}, \texttt{promises.md},
\texttt{promise\_estimation.md}.

%-----------------------------------------------------------------------------
\section{The target: an undirected model on tiles}

The tile set defines a Boltzmann distribution on tile grids $x$,
\[
  p(x) \propto e^{-E_p(x)/T}, \qquad E_p = \textstyle\sum \text{pair tables} + \sum \text{unaries},
\]
and a guarantee $G$ (connectivity, flow, support) restricts it to
$p_G(x) \propto p(x)\,\1[G(x)]$.

\paragraph{Promises are deterministic functions of the tiles.} For every block $b$ of side
$h \ge K$ in a fixed quadtree, $\pi_b = A_b(x)$, with an exact merge
$\pi_{\text{parent}} = \mathrm{merge}(\pi_{\text{children}})$. Adding them to the target
adds only hard, undirected factors:
\[
  p_G(x,\pi) \propto p(x)\,\1[G(x)] \prod_b \1[\pi_b = A_b(x)] ,
\]
whose $x$-marginal is still $p_G$. \emph{Specifying} promise values $\pi^*$ on a set of
blocks $B$ is ordinary conditioning, $p_G(x \mid \pi_B = \pi^*)$. The induced distribution
of promise pyramids is the push-forward $q = A_\# p_G$.

The target has no levels and no direction; everything below is about sampling it with local
operations.

%-----------------------------------------------------------------------------
\section{The generator: a chain graph}

Levels are indexed by cell side $h \in \{m, m/2, \dots, 1\}$ ($m$ the exemplar side).
Level $h$ is one \emph{chain component} $C_h = \{S_h, \pi_h\}$:
\begin{itemize}
  \item $S_h : \Z_{n/h}^2 \to \Z_m^2$, exemplar coordinates (an auxiliary latent);
  \item $\pi_h$, the promise grid, present only for $K \le h \le h_{\text{top}}$.
\end{itemize}
The tiles form a last component. The generator is a product of \emph{locally normalised}
level conditionals, directed coarse to fine:
\begin{equation}
  g(S,\pi,x) = g_{\text{top}}(C_{\text{top}}) \prod_{h} g_h(C_h \mid C_{2h}) \; g(x \mid S_1, \pi_K),
\end{equation}
\begin{equation}
  g_h(C_h \mid C_{2h}) = \frac{1}{Z_h(C_{2h})}\,
  \underbrace{\psi_h(S_h,S_{2h})\,\psi^\pi_h(\pi_h,\pi_{2h})}_{\text{directed, between levels}}\;
  \underbrace{\varphi_h(S_h)\,\varphi^\pi_h(\pi_h)\,\chi_h(S_h,\pi_h)}_{\text{undirected, within level}} ,
\end{equation}
\begin{equation}
  g(x \mid S_1, \pi_K) \propto \tau(x, S_1)\, e^{-E_p(x)/T} \prod_{b\,:\,h_b=K} \1[A_b(x) = \pi_b].
\end{equation}

\begin{center}
\begin{tabular}{@{}lll@{}}
\toprule
Factor & Form & Code \\
\midrule
$\psi_h$ & $k_h(S_h - U S_{2h})$, prolongation $U$ plus a noise kernel & \texttt{CoordVar.init\_children} \\
$\varphi_h$ & $e^{-\mathcal{E}_h(S_h)/T}$, patch energy against the exemplar & \texttt{CoordVar.energy} \\
$\psi^\pi_h$ & $\1[\mathrm{merge}(\pi_h)=\pi_{2h}]\,e^{-c_h(\pi_h,\pi_{2h})}$ & \texttt{refine}, \texttt{Tables.parent} \\
$\varphi^\pi_h$ & $\1[\mathrm{consistent}(\pi_h)]\,e^{-\sum u_h - \sum e_h}$ & \texttt{consistent}, \texttt{Tables.u/pair} \\
$\chi_h$ & $e^{-\lambda\,\mathrm{cost}_h(S_h,\pi_h)}$, texture honours promise & \texttt{CoordPromiseCoupling} \\
top & $\varphi_{\text{top}}\,\varphi^\pi_{\text{top}}\,\1[G(\pi_{\text{top}})]\,\chi_{\text{top}}$ & \texttt{top} \\
$\tau$ & coupling of tiles to texture, $e^{-w\,d(x, E[S_1])}$ & \texttt{heights} stage \\
\bottomrule
\end{tabular}
\end{center}

\paragraph{Why a chain graph and not one undirected energy.} With a single global
normaliser the fine levels would feed back on the coarse ones: the effective coarse
distribution becomes $\varphi_{2h}\,Z_h(C_{2h})$, tilted by how many fine configurations
each coarse value admits. Local normalisation buys three things the design depends on:
\begin{enumerate}
  \item \textbf{Locality.} A window of level $h$ depends only on a padded window of level
        $2h$, so chunked, infinite generation is well defined.
  \item \textbf{Meaningful tables.} A level's factors can be fitted as a conditional given
        the parent (counts now, pseudo-likelihood in \texttt{plfit.py}), and mean what they
        were fitted to mean.
  \item \textbf{Ancestral sampling.} Sampling $C_{\text{top}}$, then each $C_h$ given
        $C_{2h}$, is exact for the generator (Section~\ref{sec:sampler}).
\end{enumerate}

\paragraph{Promises point both ways.} In the target, promises are defined bottom-up
($x \to \pi_K \to \pi_{2K} \to \cdots$); in the generator they are sampled top-down. The two
directions cannot share one chain graph: with $\chi_h$ linking $S_h$ and $\pi_h$,
$S_h \to \cdots \to S_1 \to x \to \pi_K \to \cdots \to \pi_h - S_h$ would be a partially
directed cycle. The bottom-up definition is therefore not an edge of the generator but a
condition on it (Section~\ref{sec:agree}).

%-----------------------------------------------------------------------------
\section{When the generator equals the target}
\label{sec:agree}

Let $\hat q$ be the generator's marginal over $\pi$ and $g(x \mid \pi)$ its conditional over
tiles with $S$ marginalised out. Then
\begin{gather*}
  \hat q = q = A_\# p_G
  \quad\text{and}\quad
  g(x \mid \pi) = p_G(x \mid A(x) = \pi) \\
  \Longrightarrow\quad
  g(x) = \textstyle\sum_\pi q(\pi)\, p_G(x \mid A(x)=\pi) = p_G(x).
\end{gather*}
Neither condition is needed for $G$: every consistent $\pi$ with $G(\pi_{\text{top}})$
implies $G$ for any $x$ that honours it (the soundness theorem of each promise language).
The conditions measure how \emph{typical} the output is, not whether it is valid.

Three known departures:
\begin{itemize}
  \item $\varphi_h$ encodes the exemplar's statistics, which match $p$ locally but not at
        macro scale. $S$ is there to cross barriers local Gibbs on $p$ cannot; the
        relaxation on $E_p$ afterwards pulls $x$ back toward $p$.
  \item $\hat q$ is fitted to a corpus. It approaches $q$ only if the corpus is drawn from
        $p_G$ (e.g.\ by an exact reference sampler, not as repaired samples $R_\# p$) and the
        fitted model generalises to the configurations the sampler visits.
  \item A contract weaker than $A_b(x) = \pi_b$ (e.g.\ promises as lower bounds) preserves
        $G$ but not the second condition. The languages of Section~\ref{sec:machinery} use
        equality.
\end{itemize}

\paragraph{How far off, exactly.} Suppose the tile level samples its conditional exactly,
$g(x \mid \pi_K) = p_G(x \mid A_K(x) = \pi_K)$. Then $g(x) = \sum_\pi \hat q(\pi)\, p_G(x \mid
\pi)$, and since $\pi_K = A_K(x)$ is a function of $x$,
\[
  \mathrm{TV}(g, p_G) = \mathrm{TV}(\hat q_K, q_K).
\]
Every error of the fitted promise distribution at level $K$ passes to the output, and nothing
else does. More sweeps do not reduce it: sweeps equilibrate a level given the level above, and
nothing below ever pushes back on a coarser level.

\paragraph{Conditioning.} Clamping $\pi_B = \pi^*$ in the generator reproduces the target's
conditional only if $\hat q(\pi \mid \pi_B = \pi^*) = q(\pi \mid \pi_B = \pi^*)$. Upward
this is deterministic (merge fixes the ancestors that $B$ fully covers). Downward, a
top-down sampler must avoid coarse values with no consistent completion that reaches
$\pi^*$, which needs a backward feasibility message
\[
  \beta_h(\pi_h) = \1\big[\pi_h \text{ has a consistent refinement compatible with } \pi^*\big],
\]
as in forward-filtering backward-sampling.

%-----------------------------------------------------------------------------
\section{The sampler: prolong, perturb, smooth}
\label{sec:sampler}

Each level conditional $g_h(C_h \mid C_{2h}) \propto \psi\,\varphi$ is sampled by splitting
it into its two kinds of factor:
\begin{center}
\begin{tabular}{@{}lll@{}}
\toprule
Step & Operation & Role \\
\midrule
Prolong & $\hat C_h = $ location of $\psi_h,\psi^\pi_h$ given $C_{2h}$ & mode of the directed factor \\
Perturb & $C_h^{(0)} \sim \psi_h(\,\cdot\,, C_{2h})$ & exact draw from the directed factor \\
Smooth  & $C_h = R_h^{\,n_h}(C_h^{(0)})$ & short-run MCMC/ICM on the within-level factors \\
\bottomrule
\end{tabular}
\end{center}
Prolong + perturb is ancestral sampling of the directed part; smoothing is truncated
(block) Gibbs on the undirected part. The full pass,
\[
  C_1 = R_1^{n_1} \circ K_1 \circ R_2^{n_2} \circ K_2 \circ \cdots \circ R_{\text{top}}(\varnothing),
  \qquad K_h = \text{perturb} \circ \text{prolong},
\]
is cascadic multigrid (prolong and smooth, never restrict), or multiscale ancestral sampling
with short-run MCMC correction. The smoother is exact for $g_h$ only if its target keeps
the directed factor; otherwise the parent survives only through the initialisation and the
small $n_h$.

\begin{center}
\begin{tabular}{@{}llll@{}}
\toprule
Variable & Prolong & Perturb & Smooth \\
\midrule
$S_h$ & $U S_{2h}$ & jitter $k_h$ & ICM/Gibbs on $\varphi_h\chi_h$ ($\psi_h$ dropped) \\
$\pi_h$ & MAP over $F$ & draw over $F$ & exact chain-block Gibbs on \\
 & (Sec.~\ref{sec:machinery}) & & $\psi^\pi_h\varphi^\pi_h\chi_h$ ($\psi^\pi_h$ kept) \\
$x$ & feasible $x$ nearest & none & MCMC on $\tau e^{-E_p/T}$ that keeps \\
 & $E[S_1]$ & & $A_b(x) = \pi_b$ exactly \\
\bottomrule
\end{tabular}
\end{center}

%-----------------------------------------------------------------------------
\section{Promise machinery}
\label{sec:machinery}

This section states what a promise language must provide and how the sampler uses it, for any
guarantee $G$. Section~\ref{sec:support} works it through for support.

\subsection{What a language provides}

A block $b$ of side $h$ has an exact state $s_b$, computed from $x_b$, with an exact merge
$s_{\text{parent}} = M(s_{\text{children}})$. A promise language supplies:
\begin{itemize}
  \item an \textbf{abstraction} $A$ from states to values, $\pi_b = A(s_b)$, depending only on
        the block's own contents, so the language is translation invariant;
  \item a \textbf{promise merge} $m$ that commutes with it:
        $A(M(s_1,\dots,s_4)) = m(A(s_1),\dots,A(s_4))$;
  \item \textbf{validity} $\mathcal V_{\mathrm{pos}}(\pi)$: which values a block may take at
        its position (boundary conditions);
  \item a \textbf{seam relation} $\pi \sim_d \pi'$ for blocks adjacent in direction $d$;
  \item the \textbf{contract} at level $K$: $x$ honours $\pi_K$ iff $G(x)$ and
        $A_b(x) = \pi_b$ for every $K$-block.
\end{itemize}
A grid of values is \emph{consistent} if every value is valid, every seam obeys $\sim$, and
every $2\times2$ group merges to its parent. The language must have four properties.
\begin{enumerate}
  \item \textbf{Soundness.} A consistent pyramid whose $K$-blocks are honoured implies $G$.
  \item \textbf{Completeness.} $G(x)$ implies that $A(x)$ is a consistent pyramid. Then $q$
        lives on consistent pyramids, and a corpus can be abstracted directly, with no
        projection.
  \item \textbf{Realisability.} Every consistent value has a consistent refinement for any
        consistent halo, and every consistent level-$K$ grid has a completion $x$. Top-down
        sampling then never gets stuck.
  \item \textbf{Exact seam rule.} $\pi \sim_d \pi'$ iff some pair of states with those
        abstractions obeys the tile-level rule across the seam. A stronger rule is sound but
        forbids valid configurations, which shows up as systematic artifacts. A weaker rule is
        unsound.
\end{enumerate}

\subsection{Feasible sets and chain blocks}

For a parent $p$ at level $2h$, the \emph{feasible set} is
\[
  F(p) = R\big(\pi_{2h}(p) \mid \text{halo}\big) \cap \beta_h \cap \text{clamps},
\]
where $R$ is the set of consistent refinements given the neighbours' current children, and
$\beta_h$ is the conditioning message below. Every step for $\pi_h$ (prolong, perturb and
smooth) chooses within $F$, so every level is consistent after every step.

\paragraph{Seam interfaces.} The seam relation reads only a projection $\iota_d(\pi)$, the
\emph{interface} a value presents across side $d$. A parent's value need not fix what its
children present at the parent's external seams. When it doesn't, refinement depends on the
halo, and the level conditional couples neighbouring parents. A per-parent draw is then exact
only given the halo, which is fine for smoothing but not for the first draw of a level.

\paragraph{Chain blocks.} Choose update blocks whose internal hard constraints form a chain
$c_1,\dots,c_L$ (a path of cells or groups). Given everything outside the block, its
conditional has the form
\[
  p(c_{1:L} \mid \text{rest}) \propto
  \exp\!\Big(-\tfrac1T\Big[\sum_i e_i(c_i) + \sum_i t_i(c_i, c_{i+1})\Big]\Big),
\]
where $e_i$ collects everything fixed during the draw (unary table, parent term, pair terms
with the flanking blocks, $\chi$, clamps and $\beta$ as $+\infty$), and $t_i$ holds the hard
merge and seam masks plus soft vertical pair terms. It is drawn exactly by forward-filtering
backward-sampling:
\[
  \mu_L = 0, \qquad
  \mu_i(c) = -T\log\sum_{c'} e^{-[t_i(c,c') + e_{i+1}(c') + \mu_{i+1}(c')]/T},
\]
then $c_1 \propto e^{-[e_1 + \mu_1]/T}$ and $c_{i+1} \propto
e^{-[t_i(c_i,\cdot) + e_{i+1} + \mu_{i+1}]/T}$. Taking $T \to 0$ turns the sums into minima
and the draw into the MAP chain. So the three sampler steps are one computation:
\begin{center}
\begin{tabular}{@{}ll@{}}
\toprule
Prolong & MAP chain ($T = 0$) over $F$ \\
Perturb & exact draw ($T > 0$) over $F$ \\
Smooth  & block Gibbs, redrawing chain blocks given their neighbours, alternating block families \\
\bottomrule
\end{tabular}
\end{center}
When the transitions depend on $c_i$ only through its interface, the messages live on interface
values, not on whole values. Small interfaces are therefore a design goal for a language.

\subsection{Forward-filtering backward-sampling}
\label{sec:ffbs}

\paragraph{Problem.} A chain has $L$ sites, each taking one of $V$ values. It has site energies
$e_i(c)$ for $i = 1,\dots,L$, transition energies $t_i(c, c')$ between sites $i$ and $i+1$, and
a temperature $T \ge 0$. A hard constraint is an energy of $+\infty$. The target distribution
is
\[
  p(c_{1:L}) = \frac{1}{Z}\exp\!\Big(-\frac{E(c_{1:L})}{T}\Big), \qquad
  E(c_{1:L}) = \sum_{i=1}^{L} e_i(c_i) + \sum_{i=1}^{L-1} t_i(c_i, c_{i+1}).
\]
Define the \emph{soft minimum} at temperature $T$,
\[
  \mathrm{smin}_T\, x = -T \log \textstyle\sum_{c} e^{-x(c)/T}, \qquad \mathrm{smin}_0\, x = \min_c x(c),
\]
with $\mathrm{smin}_T$ of an all-$+\infty$ vector equal to $+\infty$.

\paragraph{Algorithm.} Algorithm~\ref{alg:ffbs} computes, for each site, the energy of the best
(or, at $T > 0$, the softened sum over all) completions of the sites after it, and then samples
from the first site forward. This is the textbook algorithm with the chain's direction
reversed: the messages run from the end, and sampling runs from the start. The code
(\texttt{envelope\_var.ffbs}) vectorises it over many independent chains.

\begin{algorithm}[ht]
\caption{Forward-filtering backward-sampling on a chain (messages from the end, samples from the start).}
\label{alg:ffbs}
\begin{algorithmic}[1]
\Require site energies $e_{1:L}$, transition energies $t_{1:L-1}$, temperature $T$, noise
\State $\mu_L(c) \gets 0$ for all $c$ \Comment{cost-to-go after the last site}
\For{$i = L-1, \dots, 1$}
  \State $\mu_i(c) \gets \mathrm{smin}_T\big[\, t_i(c, \cdot) + e_{i+1}(\cdot) + \mu_{i+1}(\cdot) \,\big]$ for all $c$
\EndFor
\State $a_1(c) \gets e_1(c) + \mu_1(c)$
\If{$a_1(c) = +\infty$ for every $c$} \State \textbf{raise} infeasible \EndIf
\State $c_1 \gets \textsc{Pick}_T(a_1)$
\For{$i = 2, \dots, L$}
  \State $a_i(c) \gets t_{i-1}(c_{i-1}, c) + e_i(c) + \mu_i(c)$
  \State $c_i \gets \textsc{Pick}_T(a_i)$
\EndFor
\State \Return $c_{1:L}$
\Statex
\Function{Pick$_T$}{$a$}
  \If{$T = 0$} \State \Return $\arg\min_c a(c)$ \Comment{ties broken by noise}
  \Else \State draw $g(c) \sim \mathrm{Gumbel}(0,1)$ independently
        \State \Return $\arg\min_c \big[a(c)/T - g(c)\big]$ \Comment{a draw with $p(c) \propto e^{-a(c)/T}$}
  \EndIf
\EndFunction
\end{algorithmic}
\end{algorithm}

\paragraph{Why it is exact.} By induction from the end of the chain,
\[
  e^{-\mu_i(c)/T} = \sum_{c_{i+1:L}} \exp\!\Big(-\tfrac1T\Big[t_i(c, c_{i+1})
    + \textstyle\sum_{j > i} e_j(c_j) + \sum_{j > i} t_j(c_j, c_{j+1})\Big]\Big),
\]
the total weight of all completions after site $i$ given $c_i = c$. So
\[
  p(c_1) \propto e^{-a_1(c_1)/T}, \qquad
  p(c_i \mid c_{1:i-1}) = p(c_i \mid c_{i-1}) \propto e^{-a_i(c_i)/T},
\]
and the product of these conditionals is $p(c_{1:L})$. So one pass returns an exact sample.
\begin{itemize}
  \item At $T = 0$ the sums become minima, and the pass returns a lowest-energy chain (the MAP
        chain, the Viterbi algorithm).
  \item Every value chosen has finite $a_i$. Every such value has at least one finite
        completion, since $\mu_i < \infty$. So once $a_1$ has a finite entry, no dead end can
        occur later.
\end{itemize}

\paragraph{Normaliser and likelihood.} The same messages give the partition function,
\[
  -T \log Z = \mathrm{smin}_T\, a_1,
\]
so the exact log-probability of an observed chain $c_{1:L}$ is
\[
  \log p(c_{1:L}) = -\frac{E(c_{1:L})}{T} - \log Z.
\]
This is the chain-block likelihood used to fit the level factors
(\texttt{envpredict.column\_nll}).

\paragraph{Cost.} Each message is a $V \times V$ reduction, so a chain costs $O(L V^2)$ time and
$O(L V)$ memory. Implementations work in log space: $\mathrm{smin}$ subtracts the minimum
before exponentiating, and $+\infty$ entries contribute nothing.

\paragraph{How the sampler uses it.} A chain is one column of promise blocks at a level.
\begin{itemize}
  \item $e_i$ holds the unary and parent tables, the pair terms with the fixed neighbouring
        columns, $\chi$, and $+\infty$ outside the feasible sets.
  \item $t_i$ holds the soft vertical pair term, plus $+\infty$ for pairs the merge or seam
        rule forbids.
  \item Prolong is the $T = 0$ call, perturb a $T > 0$ call, and each smoothing step a $T > 0$
        call on every chain of one parity given the other.
\end{itemize}

\subsection{Conditioning: precomputing feasibility}

A request is a set of clamps $\mathcal C = \{(h, b) \mapsto \pi^*_b\}$ at any levels $h \ge K$.
Before any sampling, two structures are computed bottom-up from $K$ to $h_{\text{top}}$.
\begin{itemize}
  \item \textbf{Feasible values} $F_h(b)$: the values block $b$ may take and still admit a
        consistent completion that honours every clamp below it.
  \item \textbf{Feasible seams} $B_h(b, b') \subseteq {\sim}$ for each pair of blocks adjacent
        across a hard seam: the value pairs that admit such completions jointly.
\end{itemize}
At level $K$, and at any level for a block with nothing clamped below it,
\[
  F_K(b) = \begin{cases} \{\pi^*_b\} & (K, b) \in \mathcal C \\
                         \{\pi : \mathcal V_{\mathrm{pos}(b)}(\pi)\} & \text{otherwise,} \end{cases}
  \qquad
  B_K(b, b') = {\sim} \cap \big(F_K(b) \times F_K(b')\big).
\]
One level up, write $\mathrm{ch}(p)$ for the children of $p$, and
$\mathrm{Real}_h(P, p)$ for the refinements of value $P$ that stay inside the feasible sets:
\[
  \mathrm{Real}_h(P, p) = \big\{c \in \textstyle\prod_{i \in \mathrm{ch}(p)} F_h(i) :
    m(c) = P,\ \text{every internal seam pair of } c \text{ in } B_h \big\}.
\]
Then
\[
  F_{2h}(p) = \big\{P : \mathrm{Real}_h(P, p) \neq \varnothing\big\} \cap \text{clamps at } (2h, p),
\]
and $(P, P') \in B_{2h}(p, p')$ iff there are $c \in \mathrm{Real}_h(P, p)$ and
$c' \in \mathrm{Real}_h(P', p')$ whose facing children are pairwise in $B_h$.

\paragraph{Arc consistency.} A value that has no partner in some neighbour's feasible set, under
$B$, can never be completed. So each level is pruned to a fixed point, sweeping along each chain
in both directions:
\[
  F(b) \leftarrow \{\pi \in F(b) : \exists \pi' \in F(b')\ \text{with}\ (\pi, \pi') \in B(b, b')\}
  \quad \text{for every neighbour } b',
\]
with $B$ restricted to the pruned sets.

\paragraph{Failure.} If any $F$ becomes empty, the clamps are infeasible. This is reported
before sampling starts, never discovered as a dead end in the middle of a pass.

\paragraph{Why seam sets are needed.} Per-block sets alone are not enough. The coarse seam rule
can admit two neighbouring parents whose feasible refinements disagree across their shared seam,
for example when clamps fix which half of each parent holds the matter. Arc consistency on
$B_h$ catches this. It is still only pairwise, so a combination of three or more blocks could in
principle be infeasible. In practice (Section~\ref{sec:support}) no dead end appeared in 1500
random clamp patterns.

\paragraph{What $\beta$ is.} $F_h$ and $B_h$ together are the backward message $\beta_h$ of
forward-filtering backward-sampling, carried across levels. But they are its $0/1$ support
only, not its weight. That difference shapes the distribution the pass samples, as described
below.

\subsection{The tile level}

The tile conditional is
\[
  g(x \mid S_1, \pi_K) \propto \tau(x, S_1)\, e^{-E_p(x)/T}\, \1[G(x)]
  \prod_{b\,:\,h_b = K} \1[A_b(x) = \pi_b].
\]
Split $x$ into the part $y = \rho(x)$ that the guarantee reads (e.g.\ solidity) and the rest,
and reduce $\pi_K$ to a feasible set $Y(\pi_K)$ for $y$.
\begin{itemize}
  \item \textbf{Prolong}: $y^0 = \arg\min_{y \in Y(\pi_K)} d\big(y, \rho(E[S_1])\big)$, a
        projection used only as the starting point.
  \item \textbf{Smooth}: MCMC on the conditional whose moves stay inside $Y(\pi_K)$.
        \begin{itemize}
          \item Local moves are checked locally, since the contract is per block plus seams.
          \item Local moves alone can freeze, because a block's abstraction can pin which cells
                carry a property. Moves that leave $A$ invariant but relocate that role (e.g.\
                permutations within a block that preserve its value) make the chain
                irreducible on $Y(\pi_K)$.
          \item The rest of $x$ is resampled within its class.
        \end{itemize}
\end{itemize}
The projection's bias decays with sweeps, and the limit is the exact conditional. Earlier
pipelines instead used the projection as the output and repaired around it.

\subsection{A conditioned pass, step by step}

Given clamps $\mathcal C$, a seed and the fitted level factors:
\begin{enumerate}
  \item \textbf{Precompute} $F_h$ and $B_h$ for every promise level (previous subsection).
        Stop here if any set is empty.
  \item \textbf{Top level} $h_{\text{top}}$. The chain blocks have unaries
        \[
          e_i(c) = u_{h}(c) + \lambda\,\mathrm{cost}_h\big(S_h(b_i), c\big)
                   + \infty \cdot \1[c \notin F_h(b_i)],
        \]
        and transitions
        \[
          t_i(c, c') = v_h(c, c') + \infty \cdot \1[(c, c') \notin B_h(b_i, b_{i+1})],
        \]
        where $u_h$ is the unary term, $v_h$ the soft pair term along the chain, and the cost term
        is $\chi_h$.
        \begin{itemize}
          \item Draw the even blocks first, ignoring their pair terms to other blocks.
          \item Draw the odd blocks given both even neighbours; the soft pair terms across become
                unaries.
          \item Then run $n_{\text{top}}$ sweeps of block Gibbs, alternating even and odd.
        \end{itemize}
  \item \textbf{Each level $h = h_{\text{top}}/2, \dots, K$}, given the finished level $2h$:
        \begin{enumerate}
          \item \emph{Coordinates}: prolong $S_h = U S_{2h}$ and perturb with $k_h$.
          \item \emph{Promises, prolong or perturb}: the same chain draw as above, over child
                chains, with the parent term added to the unary and the merge added to the
                transitions:
                \begin{align*}
                  e_i(c) &= u_h(c) + c_h\big(c, \pi_{2h}(\mathrm{par}(b_i)), \mathrm{pos}(b_i)\big)
                           + \lambda\,\mathrm{cost}_h\big(S_h(b_i), c\big)
                           + \infty \cdot \1[c \notin F_h(b_i)], \\
                  t_i(c, c') &= v_h(c, c') + \infty \cdot \1[\text{pair not allowed}],
                \end{align*}
                where $\mathrm{par}(b)$ is the parent of block $b$, $\mathrm{pos}(b)$ its position
                within that parent, and the parent term $c_h$ comes from $\psi^\pi_h$. A pair is
                allowed when:
                \begin{itemize}
                  \item for siblings under one parent, it is a realising pair of that parent's
                        value, $(c, c') \in \mathrm{Real}_h(\pi_{2h}(\mathrm{par}), \cdot)$
                        restricted to the chain;
                  \item for a pair across two parents, $(c, c') \in B_h$.
                \end{itemize}
                $T = 0$ gives the MAP chain (prolong); $T > 0$ gives an exact draw (perturb).
          \item \emph{Smooth}: $n_h$ sweeps. Each sweep redraws the even chains given the odd
                ones and then the odd given the even, each an exact block-Gibbs step. It then
                runs the coordinate smoother on $\varphi_h\chi_h$ against the current promises.
        \end{enumerate}
  \item \textbf{Below $K$}: coordinates only. $\chi$ couples each cell to the promise of its
        level-$K$ ancestor block.
  \item \textbf{Tile level}: reduce $\pi_K$ to $Y(\pi_K)$, project $E[S_1]$ onto it, run the
        constrained MCMC, then resample materials within their classes.
\end{enumerate}

\paragraph{Invariants.} After every step of step 3, level $h$:
\begin{itemize}
  \item is consistent with itself and with level $2h$;
  \item equals the clamps at level $h$;
  \item has every value inside $F_h$, so a consistent completion honouring every finer clamp
        exists, up to the pairwise caveat above.
\end{itemize}
The tile level honours $\pi_K$ exactly, and merges are exact. So every clamp, at every level,
holds in the output $x$, together with $G(x)$.

\paragraph{Which distribution this samples.} Let every chain draw be exact, with infinitely many
sweeps. Then the pass samples the generator with each level conditional renormalised over its
feasible set:
\[
  \tilde g(\pi) = g_{\text{top}}^{F}(\pi_{\text{top}})
  \prod_h \frac{g_h(\pi_h \mid \pi_{2h})\,\1[\pi_h \in F_h]}{Z^F_h(\pi_{2h})},
  \qquad
  Z^F_h(\pi_{2h}) = \sum_{\pi_h \in F_h} g_h(\pi_h \mid \pi_{2h}).
\]
The exact conditional $g(\pi \mid \mathcal C)$ instead weights each coarse value by the
probability mass of its completions that reach the clamps. That is the soft backward message
\[
  \beta^{\text{soft}}_h(\pi_h) = \Pr_g\big(\text{finer levels honour } \mathcal C \mid \pi_h\big),
\]
whose support is $F_h$. The two agree when $\beta^{\text{soft}}$ is constant on its support.
Otherwise the coarse levels of $\tilde g$ over-weight values with few feasible completions.
The soft message is a sum over all finer pyramids, which is intractable in two dimensions. Along a
single chain it is a forward pass, so chain-structured approximations are possible. This is a
further departure to add to those in Section~\ref{sec:agree}: with clamps,
\[
  \mathrm{TV}\big(g(\cdot \mid \mathcal C),\, p_G(\cdot \mid \mathcal C)\big)
\]
depends both on $\hat q$ and on how far $\tilde g$ is from $g(\cdot \mid \mathcal C)$.

\subsection{Fitting the level factors}

A table indexes whole values, so a value never seen whole in the corpus gets an uninformed
energy. A linear predictor over language features generalises from the parts of a value that
were seen:
\[
  E_\theta(\pi_h \mid \text{context}) = \theta_h \cdot \phi(\pi_h, \text{context}),
\]
with $\phi$ built from the value's components, its interfaces and its relation to the parent.
Tables are the case where $\phi$ is one-hot on whole values. Fit $\theta_h$ by \emph{chain-block
pseudo-likelihood}, maximising $\sum_{\text{blocks}} \log p_\theta(c_B \mid \text{rest})$ over the corpus.
Each term is exact by the forward pass above, and it is the conditional the smoother draws from.

Two requirements come from the data rather than the model:
\begin{itemize}
  \item \textbf{Coverage.} The corpus must show the configurations the sampler will visit,
        relative to the block grid: surfaces at every offset within a block, structures
        crossing seams, and so on. Outside its data, a model extrapolates by its inductive
        bias, which can be confidently wrong.
  \item \textbf{Agreement of inputs.} The corpus (which sets $\hat q$), the exemplar (which
        sets $\chi$ and $\tau$) and the target $p$ (the tile level) must describe the same
        world. Where they disagree, the terms fight, and the stronger one wins locally.
\end{itemize}

%-----------------------------------------------------------------------------
\section{Example: support as local height envelopes}
\label{sec:support}

Code: \texttt{quantities/envelope.py} (language), \texttt{envelope\_var.py} (promise levels),
\texttt{heights.py} (tile level), \texttt{refheights.py} (exact reference sampler),
\texttt{envpredict.py} (linear predictor); experiments in \texttt{notes/experiments/}.

\paragraph{Guarantee.} Side view with gravity. The bottom $G$ rows are a ground band. A solid
cell outside the band needs a solid or band cell directly below. Walking down from any
column's topmost solid shows that the valid maps are exactly the heightmaps: column $c$ is air
above a height $H(c) \ge G$ and solid below it.

\paragraph{State.} For a $k \times k$ block: the local height $t(c) \in [0,k]$ of each column
(the solid run from the block's bottom row, so $t = \mathrm{clip}(H - \text{offset}, 0, k)$),
the top-row solid mask, and a count of violations. With upper child $u$ over lower child $l$,
the merge is
\[
  t = \begin{cases} t_l & t_l < k \\ k + t_u & t_l = k \end{cases},
  \qquad \text{viol} = \textstyle\sum \text{viol} + \#\{t_u > 0 \wedge \neg\,\mathrm{top}_l\}.
\]

\paragraph{Values.} Each block side is split into $I$ column intervals. For each interval the
value records the bins of $\min t$ and $\max t$, where
\[
  b(t) = 0 \text{ if } t = 0, \qquad b(t) = F = Q+1 \text{ if } t = k, \qquad
  b(t) = \lceil tQ/k \rceil \text{ otherwise}.
\]
With $C = Q+2$ bins there are $C(C+1)/2$ ordered pairs per interval and $V = (C(C+1)/2)^I$
values. $I = Q = 2$ gives $V = 100$.

\paragraph{Merge.} Stacked children map to the parent's coarser bins by
$\mathrm{up}(0) = 0$, $\mathrm{up}(b) = \lceil b/2 \rceil$ for interior $b$,
$\mathrm{up}(F) = Q/2$, and $\mathrm{up}^+(b) = Q/2 + \mathrm{up}(b)$ with $\mathrm{up}^+(F) = F$
for the upper child:
\[
  \mathrm{lo} = \begin{cases} \mathrm{up}(l.\mathrm{lo}) & l.\mathrm{lo} < F \\
                              \mathrm{up}^+(u.\mathrm{lo}) & \text{else} \end{cases},
  \qquad
  \mathrm{hi} = \begin{cases} \mathrm{up}^+(u.\mathrm{hi}) & u.\mathrm{hi} > 0 \\
                              \mathrm{up}(l.\mathrm{hi}) & \text{else} \end{cases}.
\]
Side by side, lo is the minimum and hi the maximum of the two child intervals a parent interval
covers. The merge is exact because a column that is not full in the lower child must be empty
in the upper one. Nesting bins with $Q$ even makes the bin arithmetic exact as well.

\paragraph{Seam rule, validity, interfaces.} For an upper value over a lower one, per interval:
\[
  u.\mathrm{hi} > 0 \Rightarrow l.\mathrm{hi} = F, \qquad
  u.\mathrm{lo} > 0 \Rightarrow l.\mathrm{lo} = F.
\]
If any upper column has matter at its bottom row, some lower column is full; if all do, all
are. This is exactly the image of the per-column rule. The rule it replaced, ``any matter above
$\Rightarrow$ all columns full'', was stronger, and it produced slab artifacts.
\begin{itemize}
  \item Side-by-side seams carry no hard constraint. The torus wrap seam, from the band up to
        the top row, is exempt.
  \item Validity: a ground-row block has $\mathrm{lo} \ge b(G)$.
  \item Interfaces: an upper value presents $(\mathrm{lo}>0, \mathrm{hi}>0)$ and a lower value
        presents $(\mathrm{lo}=F, \mathrm{hi}=F)$ per interval, a few bits in all.
\end{itemize}

\paragraph{Properties.} Completeness holds because every heightmap abstracts to a consistent
pyramid. Realisability follows from clipping being monotone: in any stack of blocks, every
block's lo is attained at the interval's lowest column and every hi at its highest. A stack is
therefore consistent iff it is the clip of a common $[H_{\min}, H_{\max}]$. Given intervals at
least two columns wide, any consistent level-$K$ grid has a completion.

\paragraph{Chain blocks.} All hard constraints run vertically, and a parent's value over child
column $d$ depends only on that child column. So a level splits into independent \emph{child
columns}: chains of $n/h$ values with the merge mask inside each parent and the seam rule
between parents. Each column is drawn exactly by forward-filtering backward-sampling with
$V = 100$ states. Horizontal pair terms become unaries when even and odd columns alternate.
The per-parent $\beta$ was observed to fail (a feasible-looking coarse seam forced an overhang
one level down); the seam relations $B_h$ fixed it, with no dead ends in 1500 random clamp
patterns.

\paragraph{A conditioned pass, worked.} Take $n = 128$, $K = 16$, $h_{\text{top}} = 64$,
$G = 8$, $I = Q = 2$. Clamp one $h = 32$ block, at block row 1 and block column 1 (rows 32--63,
columns 32--63), to \emph{full}, $\mathrm{lo} = \mathrm{hi} = F$ in both intervals: a mountain.
\begin{itemize}
  \item \textbf{Precomputation, level 16.} Nothing is clamped, so each $F_{16}(b)$ is the valid
        values; ground-row blocks need $\mathrm{lo} \ge b(8) = 1$.
  \item \textbf{Level 32.} The clamped block's set is $\{(F,F),(F,F)\}$. Its realising children
        are all full. The seam rule then forces the blocks below it in the same column to be
        full: block rows 2 and 3 at $h = 32$, which becomes columns 32--63 full up to row 64 at
        $h = 16$. So the $F$ and $B$ sets of that column shrink to ``full'' below the clamp.
  \item \textbf{Level 64.} The clamp is the lower-right child of parent $(0, 0)$. The lower child
        in that child column is full, so every column there has local height $32 + t_u \ge 32$.
        The parent's interval over those columns therefore has $\mathrm{lo},
        \mathrm{hi} \in \{1, 2, F\}$ (bin 1 is exactly height 32). By the seam rule, ground block
        $(1, 0)$ below it must be full over those columns.
  \item \textbf{Sampling.} At $h = 64$ the chain for parent column 0 draws values from these
        shrunken sets. At $h = 32$ the clamped block is fixed, and its column's blocks below are
        forced full. The other columns, and the upper-left child of $(0, 0)$, are free, subject to
        the merge with their parent. The same holds down to $h = 16$.
  \item \textbf{Tile level.} For columns 32--63 the interval constraints give $H \ge 96$. Every
        other interval is constrained only by its sampled envelope.
\end{itemize}
The output always honours the clamp. How the free neighbours are distributed depends on $\hat
q$ and on the hard-versus-soft $\beta$ effect: in the experiments, count tables fitted to a flat
corpus put a spurious shoulder beside the mountain, which the true conditional does not have.

\paragraph{Tile level.} Here $y = H$. Because $t$ is monotone in $H$, $\pi_K$ reduces to four
numbers per global interval $J$: all $H \in [a_0, b_1]$, $\min_J H \in [a_0, a_1]$ and
$\max_J H \in [b_0, b_1]$.
\begin{itemize}
  \item \textbf{Prolong}: a dynamic program over each interval, with four states for whether
        the low and high witnesses have been seen, minimising disagreement with the texture.
  \item \textbf{Smooth}:
        \begin{itemize}
          \item $H(c) \pm 1$ moves, checked in $O(1)$ with per-interval witness counts;
          \item swaps of two columns within an interval. These keep the interval's set of
                heights, so every envelope is unchanged, and they move the witness role;
          \item heat-bath over materials within the solid and air classes.
        \end{itemize}
\end{itemize}

\paragraph{Fitting.} Features are shared across interval positions:
\begin{itemize}
  \item each interval's (lo, hi) pair;
  \item the difference between adjacent intervals, and whether one is full while the other is
        empty;
  \item the interfaces of stacked intervals;
  \item each child interval against its parent interval.
\end{itemize}
Each child column's likelihood is exact by the forward pass.

\paragraph{Results} ($128^2$, $K = 16$, $h_{\text{top}} = 64$, $G = 8$).
\begin{center}
\begin{tabular}{@{}ll@{}}
\toprule
Guarantee & 0 unsupported cells and satisfaction $1.0$ at every level, in every run \\
          & (the previous pipeline repaired $28\%$ of cells) \\
Clamps & always honoured, including out-of-distribution ones \\
Agreement & unclamped, $\tau = 0$: every statistic within $|z| \le 1$ of the exact $p_G$ \\
          & (height-histogram TV $0.06$, promise TV $0.01$) \\
Cost & promise pyramid $\approx 0.03$\,s, tile level $\approx 0.2$\,s \\
\bottomrule
\end{tabular}
\end{center}
Three lessons follow from Section~\ref{sec:machinery}.
\begin{itemize}
  \item \textbf{The target itself is flat.} At the tile set's own parameters $p_G$ is flat
        ground, so its pyramid carries almost no information. A mountain clamp then leaves
        $\hat q$ extrapolating: count tables are uninformed, and the linear predictor
        confidently spreads ridges, because it learned only ``neighbours are alike''.
  \item \textbf{Coverage fixes seam snapping.} With a corpus of rolling hills, the share of
        surface columns lying exactly on a block seam falls from $0.38$--$0.48$ to $0.10$ (the
        corpus's own is $0.06$), and the flat caps disappear.
  \item \textbf{Remaining caps come from the inputs disagreeing.} The remaining flat-topped,
        stepped shapes come from the exemplar, which is built from rectangular fills, acting
        through $\tau$ and $\chi$. With $\tau$ off, the output is rolling terrain.
\end{itemize}

%-----------------------------------------------------------------------------
\section{The coordinate factors in texsyn}

\paragraph{Prolongation $U$.} For parent cell $p$ and child offset $\delta \in \{0,1\}^2$,
$(U S_{2h})(2p+\delta) = S_{2h}(p) + \lfloor (2\delta - 1)\,h/2 \rfloor \pmod m$: siblings are
$h$ apart, centred on the parent, and read the parent's exemplar window.

\paragraph{Kernel $k_h$.} Rounded uniform of half-width $h\,r_l$ ($l = L - \log_2 h$), from
hashed noise $\xi$: $J_h(p) = \lfloor h\, r_l (2\xi_p - 1) + \tfrac12 \rfloor$. With the
default $r = (1,1,0,\dots)$, $k_h$ is a point mass below the top two levels.

\paragraph{Energy $\mathcal{E}_h$.} With $F_h = G_{h/2} * F$ the blurred one-hot
signature-and-socket features, $P_h,\mu_h$ the level's PCA basis and mean, and
$\Delta \in \{-2,\dots,2\}^2$,
\begin{align*}
  N_E^h(u) &= P_h^\top\big([F_h(u + h\Delta)]_\Delta - \mu_h\big), &
  N_S^h(p) &= P_h^\top\big([F_h(S(p+\Delta))]_\Delta - \mu_h\big),
\end{align*}
\[
  \mathcal{E}_h(S) = \sum_p \big\| N_S^h(p) - N_E^h(S(p)) \big\|^2
  \;+\; \beta\,\1[h=1]\, E_p\big(E[S]\big).
\]
This is a nonparametric, coarse-grained patch energy: the exemplar is the model at every
level, and the blur fixes which statistics a level matches.

\paragraph{Smoother $R_h$.} Four checkerboard colours; each cell $p$ of a colour picks from
the 18-element multiset
$\mathcal{C}(p) = \{S(p+\Delta) - h\Delta,\; C2_h(S(p+\Delta)) - h\Delta : \Delta \in \{-1,0,1\}^2\}$
with local energy
$e_p(c) = w(c)\,\|N_S^h(p) - N_E^h(c)\|^2 + \beta\,\1[h=1]\sum_d E^d_{\mathrm{pair}}(E[c], E[S(p+d)])$,
$w = 1$ for coherent candidates and $\kappa$ for jumps, by Gumbel-max:
\[
  P(S(p) = c) \propto \mathrm{mult}_{\mathcal{C}(p)}(c)\,e^{-e_p(c)/T},
  \qquad T = 0 \Rightarrow \text{ICM (argmin)}.
\]

\paragraph{Where $R_h$ differs from a Gibbs kernel for $g_h$.}
\begin{enumerate}
  \item \textbf{Missing parent term.} $e_p$ omits $-T\log k_h$.
  \item \textbf{Missing blanket terms.} Only $p$'s own term is scored; the 24 other
        neighbourhoods containing $S(p)$ are ignored, and $N_S^h(p)$ uses $p$'s current value.
  \item \textbf{State-dependent proposals.} $\mathcal{C}(p)$ depends on the neighbours and has
        no Metropolis--Hastings correction; multiplicities act as a $\log$-count coherence bonus.
  \item \textbf{Overlapping colours.} Same-colour cells are 2 apart, inside each other's
        $5\times5$ windows, so parallel updates read stale values (spacing $\ge 5$ fixes it).
\end{enumerate}

%-----------------------------------------------------------------------------
\begin{thebibliography}{9}
\bibitem{kf} D.~Koller and N.~Friedman, \emph{Probabilistic Graphical Models}, MIT Press,
  2009. \S4.6 (chain graphs), ch.~12 (ancestral and Gibbs sampling).
\bibitem{bishop} C.~M.~Bishop, \emph{Pattern Recognition and Machine Learning}, Springer,
  2006. \S8.1.2, \S11.3.
\bibitem{mackay} D.~J.~C.~MacKay, \emph{Information Theory, Inference, and Learning
  Algorithms}, CUP, 2003. Ch.~29--30.
\bibitem{mg} W.~L.~Briggs, V.~E.~Henson, S.~F.~McCormick, \emph{A Multigrid Tutorial},
  2nd ed., SIAM, 2000.
\bibitem{lh} S.~Lefebvre and H.~Hoppe, Parallel controllable texture synthesis,
  \emph{SIGGRAPH} 2005.
\bibitem{pl} R.~Paget and I.~D.~Longstaff, Texture synthesis via a noncausal nonparametric
  multiscale Markov random field, \emph{IEEE TIP} 7(6), 1998.
\end{thebibliography}

%-----------------------------------------------------------------------------
\section{Connectivity in a block: local auxiliary fields}

Code: \texttt{castlegen/legacy/blockfield.py}; test \texttt{notes/experiments/legacy/block\_field.py}.
Port semantics: \texttt{castlegen/legacy/blockconn.py}; the pipeline that produces the patterns:
\texttt{notes/experiments/legacy/maze.py}.

\paragraph{Notation.} All symbols of this section, defined once.
\begin{itemize}
  \item $k = 16$: block side. $P = \{0,\dots,k-1\}^2$: the block's cells, written $p, q$.
        $N(p) \subseteq P$: the in-block 4-neighbours of $p$.
  \item $\mathcal T$: the tile set (signatures), $|\mathcal T| = 56$.
        $t_p \in \mathcal T$: the tile at $p$; $x = (t_p)_{p \in P}$.
  \item $\mathrm{node}(t) \in \{0,1\}$: tile $t$ is walkable (a room, hallway or courtyard tile).
  \item $\mathrm{sock}_\sigma(t) \in \{0,1\}$: side $\sigma \in \{N,E,S,W\}$ of $t$ can join some walkable tile.
  \item $J(p,q) \in \{0,1\}$ for $q \in N(p)$: $p$ and $q$ are joined, i.e.\ both walkable and
        their facing sockets connect (the tile set's join tables).
  \item Ports: the 8 half-sides of the block, $k/2$ edge cells each. The pattern
        $\ell \in \{0, 1..4, \mathrm{OPEN}\}^8$ (from the promise, \S below) gives each port
        \emph{closed} (0), a \emph{class} (1..4: openings joined outside the block), or
        \texttt{OPEN} (must have openings, assumes nothing outside).
  \item Opening: a walkable edge cell $p$ with $\mathrm{sock}_\sigma(t_p)$ on its outward side $\sigma$.
        $\mathrm{exit}(p) \in \{0,1\}$: $p$ is an opening on a port with a class.
  \item Fields: $g_p \in \{0,\dots,G\} \cup \{\infty\}$, the number of gaps on the route $p$
        claims to an exit; $d_p \in \{0,\dots,D\}$, the route's step count (ignored when
        $g_p = \infty$). Here $G = 8$, $D = 96$. $\tilde g = g$ with $\tilde\infty = G + 1$.
  \item $s(p,q) = 1 - J(p,q)$: the gap cost of the step $p \to q$.
  \item $\lambda \ge 0$: the guidance weight, ramped $0 \to \lambda_{\max} = 8$.
        $\mu = \lambda$ (cost per gap, walkable cells), $\epsilon = 0.1\,\mu$ (walls),
        $\delta = 0.05$ (per step), $L = 1000$ (per invalid cell).
  \item $E_{\text{tile}}(x)$: the tile energy at $T = 1$: pair tables over in-block pairs,
        closed ports facing wall, $-\log z(t)$ per tile, and the parent term
        $\beta \sum_p -\log \hat p(t_p \mid W_p)$, where $W_p$ is the exemplar tiles within
        2 of $p$'s coordinate in the parent window and $\hat p$ their smoothed frequencies.
  \item $c(p) = (x+2y) \bmod 5$ for $p = (y,x)$: the update class.
\end{itemize}

\paragraph{Where the pattern comes from.} The promise levels ($h = 64, 32, 16$) give each
16-block its open sides and a key; a side is \emph{downhill} if open into a lower key. Both
halves of a downhill side get class 1, both halves of another open side \texttt{OPEN}, and
closed sides 0. Every block but the root has a downhill side.

\paragraph{Validity.} Cell $p$ is \emph{valid} when
\[
  V_p \;=\; [g_p = \infty] \;\lor\; \mathrm{exit}(p) \;\lor\; \exists\, q \in N(p):\;
  g_q < \infty,\; d_q < d_p,\; g_q + s(p,q) \le g_p .
\]
Such a $q$ is a \emph{witness} of $p$. $V_p$ reads only $p$ and $N(p)$.

\paragraph{Energy.}
\[
  E(x,g,d) = E_{\text{tile}}(x)
  + \lambda\, n_{\text{closed}}(x) + \lambda\, n_{\text{empty}}(x)
  + L \sum_p \1[\lnot V_p]
  + \sum_p \Big( \mu_p\, \tilde g_p + \delta\, d_p\, \1[g_p < \infty] \Big),
\]
with $\mu_p = \mu$ if $\mathrm{node}(t_p)$, else $\epsilon$;
$n_{\text{closed}}$ the openings on closed ports (a sum over edge cells);
$n_{\text{empty}}$ the class or \texttt{OPEN} ports with no opening.

\paragraph{What zero means.} If every cell is valid and $g_p = 0$ for every walkable $p$,
then every walkable cell is joined inside the block to an exit: follow witnesses from $p$;
$d$ strictly decreases, so the walk ends, and only at an exit; $g = 0$ forces every step to
be a join ($s = 0$), since $g_q + s \le g_p = 0$. The check is local: count invalid cells and
walkable cells with $g > 0$.

\paragraph{The update.} Every term of $E$ but $n_{\text{empty}}$ spans a cell and its
neighbours, so two cells interact only within distance 2. Cells of one class $c(p)$ are at
least 3 apart, so they share no term, and drawing them simultaneously from their single-site
conditionals is exact.

\begin{algorithm}[h]
\caption{One sweep of the field sampler}
\begin{algorithmic}[1]
\For{$c = 0, \dots, 4$}
  \ForAll{$p$ with $c(p) = c$ \textbf{in parallel}}
    \ForAll{$q \in N(p)$}
      \State $b_q \gets V_q$ with $p$ excluded as a witness
    \EndFor
    \ForAll{candidates $(t, g, d) \in \mathcal T \times (\{0..G\} \cup \{\infty\}) \times \{0..D\}$}
      \State $e \gets$ the terms of $E$ that involve $p$, with $(t_p, g_p, d_p) = (t, g, d)$:
      \State \quad tile pairs with $N(p)$, $-\log z(t)$, parent term, port terms at $p$
      \State \quad $+\,\mu_t\, \tilde g + \delta\, d\, \1[g < \infty]$
      \State \quad $+\,L\, \1[\lnot V_p]$ \Comment{own validity under $(t,g,d)$}
      \State \quad $+\,L \sum_{q \in N(p),\, \lnot b_q} \1\big[\lnot(g < \infty \land d < d_q \land g + s_t(q,p) \le g_q)\big]$
    \EndFor
    \State $(t_p, g_p, d_p) \gets \arg\max_{(t,g,d)} \big(-e(t,g,d) + \text{Gumbel}\big)$ \Comment{exact draw at $T = 1$}
  \EndFor
\EndFor
\end{algorithmic}
\end{algorithm}

\noindent Here $s_t(q,p)$ is the step cost with $p$'s tile set to $t$; $\mu_t = \mu$ if
$\mathrm{node}(t)$, else $\epsilon$. The last line of the energy charges $L$ for each
neighbour that relied on $p$ as its witness and would lose it. $n_{\text{empty}}$ couples the
$k/2$ cells of a port; the update reads the port's other cells as they are, so it alone is
not exact in parallel.

\paragraph{Initialisation.} $x$ starts as the parent window. $(g,d)$ start at the least
valid values: a Dijkstra from the exits, ordered lexicographically on $(g,d)$, with step
$q \to p$ adding $(s(p,q), 1)$; values past $G$ or $D$ become $\infty$. The start is valid,
and with $L$ large the chain stays valid.

\paragraph{Why it guides.} Walkable cells of an orphan component can only claim routes that
cross gaps, so each pays $\mu$ per gap. Walls relay routes for $\epsilon \ll \mu$. Turning a
wall next to a joined region into a joined walkable tile costs that cell nothing (it takes
$g = 0$ from its witness), and lets the cells behind it lower $g$ by one over the next
sweeps. Paths grow toward orphans as a ratchet; the reward arrives a sweep or two later,
through neighbour reads only. An invalid cell must cost a constant $L$: ramped with $\lambda$
it is free early, and one invalid cell lets a whole orphan component claim $g = 0$.

\paragraph{Closed form over $(g, d)$.} Enumerating the algorithm's candidates costs
$|\mathcal T| (G+2) (D+1) \approx 54{,}000$ per site. The implementation draws the same
conditional without enumerating. For a fixed tile $t$, the energy over the $(g,d)$ grid
($g < \infty$) is
\[
  \mu_t\, g + \delta\, d + L\Big(\1[\lnot V_p(g,d)] + \sum_{q \text{ dependent}} \1[\lnot \mathrm{sat}_q(g,d)]\Big),
\]
where own validity is a union of quadrants, $V_p = \mathrm{exit}(p) \lor \bigvee_q
\{g \ge g_q + s_q,\ d \ge d_q + 1\}$ over neighbours with $g_q < \infty$, and each dependent
($b_q$ false) is satisfied on a rectangle, $\mathrm{sat}_q = \{g \le g_q - s_q,\ d \le d_q - 1\}$.
Cutting the $g$ axis at every $g_q + s_q$ and $g_q - s_q + 1$, and the $d$ axis at every
$d_q + 1$ and $d_q$, splits the grid into at most $9 \times 9$ rectangles of constant
violation count $v$. A rectangle $[g_0, g_0 + n_g) \times [d_0, d_0 + n_d)$ has weight
\[
  e^{-Lv - \mu_t g_0 - \delta d_0}\;
  \frac{1 - e^{-\mu_t n_g}}{1 - e^{-\mu_t}}\;
  \frac{1 - e^{-\delta n_d}}{1 - e^{-\delta}}
\]
(a sum of $n$ terms when the rate is 0), and $g = \infty$ adds one entry,
$e^{-\mu_t (G+1) - L\,n_{\text{dep}}}$. The sum depends on $t$ only through its key: the gap
bits $s_q$, $\mathrm{exit}$ and $\mathrm{node}$, so each site computes it once per distinct
key. A draw picks $t$ by its tile terms plus its key's log-sum, then a rectangle (or
$\infty$) by weight, then $g$ and $d$ independently within it by inverting truncated
geometrics. Tested against enumeration (\texttt{tests/legacy/test\_blockfield.py}); 8\,ms per
sweep of a 16-block on one core, from 1.1\,s enumerated.

\paragraph{Scope.} Every class exit is a root, so the field certifies that
each walkable cell reaches \emph{some} class; required links between classes, and the root
block (no class port: one component), need more (a field per group; a designated root cell).

\end{document}
